\section{Lifting models and the proper iteration lemma}
\label{sec:p06-proper-iteration-lemma}

The following prescribed-initial master lemma uses countable-support
quotients and lifted elementary models. The related construction in
\cite[Proposition~6.4, pp.~158--159]{GuzmanHrusakTellez2020RestrictedMAD}
has additional MAD-family hypotheses and refers its final fusion back to
Abraham. None of those additional claims is adopted here. We use the
actual bounded-name, quotient, two-step and restriction constructions of
\cref{sec:p06-forcing-iterations-support} and the master criterion of
\cref{def:p06-pfm-master-condition}. Smaller conditions are stronger.

\begin{hbconvention}{def:p06-pil-datum}{The full iteration datum and hereditary height}
Let \(\mathcal I=\langle P_\xi,\dot Q_\xi:\xi<\beta\rangle\) be
the countable-support system already defined, including its raw condition
sets \(R_\xi\), quotient maps and chosen sections. Put
\(\kappa\geq\aleph_1\) above the size of its transitive closure and
take regular \(\chi>2^\kappa\). Fix countable
\(N\prec(H_\chi,\in)\) containing this full datum and \(\beta\).
For \(\alpha\in N\cap(\beta+1)\), a prescribed
\((N,P_\alpha)\)-master \(p\) need not be in \(N\).
Raw representatives of coherent initial conditions are chosen to agree
literally on their earlier domains; mutual-extension equivalence makes
such choices possible. Generic filters, when used semantically, are
supplied, not inferred from consistency or transitivity.
\end{hbconvention}

\begin{hbdefinition}{def:p06-pil-lift}{The lifted elementary model}
For a \(P\)-generic \(G\), put
\[
 N[G]=\{\tau^G:\tau\in N\text{ is a }P\text{-name}\}.
\]
It is an evaluated set of names, not the set-theoretic union \(N\cup G\).
The canonical generic name belongs to \(N\), so \(G\) itself is an
element of \(N[G]\); that does not make all members of \(G\) elements
of this countable model.
\end{hbdefinition}

\begin{hblemma}{thm:p06-pil-lifting}{proved=annex:p06.pil.lifting}{Hereditary names, elementarity and the ground trace}
If \(p\) is an \((N,P)\)-master and \(p\in G\), then
\(N[G]\prec H_\chi^{V[G]}\), it is countable, and
\[
 N[G]\cap V=N,\qquad N[G]\cap\operatorname{Ord}=N\cap\operatorname{Ord}.
\]
In particular \(N[G]\cap a=N\cap a\) for ground \(a\in N\).
The proof supplies the hereditary-small-name and witness-selection
arguments; no Power Set axiom for \(H_\chi\) is assumed.
\end{hblemma}

\begin{hbdefinition}{def:p06-pil-tail}{The ground-represented tail quotient}
After a \(P_\alpha\)-generic \(G_\alpha\), let
\[
 T_{\alpha\beta}(G_\alpha)=
 \{u\in P_\beta:u\mathbin{\upharpoonright}\alpha\in G_\alpha\}
\]
with the old full-stage order. For this poset put \(u\leq_s v\) iff
every old-order strengthening of \(u\) is compatible with \(v\),
and identify mutual \(\leq_s\). Write \(c_G(u)\) for its separative
class and \(\mathbb T_{\alpha\beta}(G)\) for the resulting tail poset.
Changing an initial prefix within \(G_\alpha\), while retaining the
same tail, gives the same separative class. This is not an assertion
that the old full conditions are literally equal or mutually extending.
\end{hbdefinition}

\begin{hblemma}{thm:p06-pil-tail-representation}{proved=annex:p06.pil.tail-representation}{Factoring and representing model tails by ground conditions}
The ground-represented quotient and its separative quotient give the
usual factorization over the initial segment. If \(p\) is an initial
master and forces \(\dot q\in\mathbb T_{\alpha\beta}(\dot G_\alpha)
\cap N[\dot G_\alpha]\), one can choose a name \(\dot\rho\),
not necessarily in \(N\), with values in \(R_\beta\cap N\), such that
\[
 p\forces [\dot\rho\mathbin{\upharpoonright}\alpha]\in\dot G_\alpha,
 \qquad c_{\dot G_\alpha}([\dot\rho])=\dot q.
\]
Every such raw value has countable support contained in \(N\cap\beta\).
Thus model tails have a uniform ground countable coordinate envelope,
without assuming that an arbitrary fresh countable support is ground.
\end{hblemma}

\begin{hblemma}{thm:p06-pil-successor-master}{proved=annex:p06.pil.successor-master}{An actually named iterand master}
If \(p\) is \((N,P)\)-generic and forces \(\dot Q\) proper and
\(\dot a\in\dot Q\cap N[\dot G]\), a bounded total condition name
\(\dot b\) can be chosen so that \(p\) forces \(\dot b\leq\dot a\)
and \(\dot b\) master for \(N[\dot G]\). Then \((p,\dot b)\) is
a master for the two-step iteration. The selected master name need not
belong to \(N\), and no closure of \(\dot Q\) is assumed.
\end{hblemma}

\begin{hbtheorem}{thm:p06-pil-extension}{proved=annex:p06.pil.extension}{The support-controlled prescribed-initial master lemma}
Assume every \(P_\xi\) forces \(\dot Q_\xi\) proper. Given the
standing \(N,\alpha,p\) and a tail name \(\dot q\) as above, there
is \(r\in P_\beta\), master for \(N\), with
\[
 r\mathbin{\upharpoonright}\alpha=p,
 \quad p\forces c_{\dot G_\alpha}(r)\leq\dot q,
 \quad \operatorname{supp}(r)\setminus\alpha\subseteq N\cap[\alpha,\beta).
\]
The support inclusion is for the constructed raw representative.
At limits, including \(\sup(N\cap\beta)<\beta\), the proof constructs
coherent initial masters and supplies the missing tail domination and
dense-set checks. It never takes a lower bound of arbitrary proper
conditions.
\end{hbtheorem}

\begin{hbcorollary}{thm:p06-pil-model-masters}{proved=annex:p06.pil.model-masters}{Master production with the iteration datum}
For every countable elementary model containing the full iteration
datum and every \(u\in P_\beta\cap N\), a master strengthening of
\(u\) exists. This is the iteration-data master-production conclusion.
Its exact model parameters are retained; no source-specific 1982
presentation or PFA/MM construction is thereby admitted.
\end{hbcorollary}

\begin{hbexample}{def:p06-pil-prefix-boundary}{Why the tail comparison is not literal full-condition extension}
Even if \(p\) forces an initial condition into the generic, it need
not literally strengthen that condition. A nonseparative poset may
have two incomparable compatible conditions whose common lower bounds
are dense below each. They have the same separative value. Prescribing
one of them as the initial condition cannot also demand literal extension
of the other. The quotient in the theorem retains this distinction.
\end{hbexample}
